\section{杂化泛函}
杂化泛函是一类特殊的交换-相关泛函，该泛函中部分地考虑了\textrm{Hartree-Fock}方法中\textrm{Fock}交换形式的贡献。

\subsection{\rm{PBE}和\rm{PBE0}得到的\ch{Si}的带隙}
分别采用\textrm{PBE}和\textrm{PBE0}泛函计算金刚石结构的\ch{Si}的带隙。

\textrm{DFT}计算中  最有效的泛函包括\textrm{LDA}、\textrm{GGA}和\textrm{meta-GGA}泛函。\textrm{GGA}和\textrm{meta-GGA}属于半局域泛函泛函，具体的泛函形式有多种多样。\textrm{VASP}中常用的\textrm{GGA}是\textrm{PBE}泛函，也是各类固体物理计算中使用得最普遍的泛函之一。有必要指出的是，\textrm{LDA}和\textrm{GGA}对于固体能带的带隙估计往往偏低，因为\textrm{Kohn-Sham}方程得到的本征态的优劣与采用的交换-相关泛函息息相关。如果采用\textrm{meta-GGA}和杂化泛函，得到的带隙往往会理想一些。

受\textrm{Hartree-Fock}方法的启发，杂化泛函是在原有的半局域泛函基础上混入了一部分\textrm{Fock}交换项的贡献，因为\textrm{Hartree-Fock}方法精确描述了体系的交换作用。对于周期体系，引入杂化泛函主要为了修正\textrm{LDA}和\textrm{GGA}泛函低估带隙的问题。非局域的\textrm{Fock}势函数可以表示为
\begin{displaymath}
	V_{\mathrm{x}}^{\mathrm{HF}}(\vec r,\vec r^{\prime})=-\dfrac{e^2}2\sum_{n\vec k}f_{n\vec k}\dfrac{\psi^{\ast}(\vec r^{\prime})\psi_{n\vec k}(|vec r)}{|\vec r-\vec r^{\prime}|}
\end{displaymath}
这里$f_{n\vec k}$是指定$\vec k$点上能带$n$上的电子占据数。$\vec r$是电子位置，$e$是电子电荷。因为\textrm{Fock}势函数依赖粒子波函数而非粒子密度，所以严格地讲，引入杂化泛函实际上破坏了密度泛函理论的基本假设。

\textrm{PBE0}是\textrm{VASP}常用的杂化泛函，其构成为
\begin{displaymath}
	E_{xc}^{\mathrm{PBE0}}=\dfrac14E_{x}^{\mathrn{HF}}+\dfrac34E_x^{\mathrm{PBE}}+E_c^{\mathrm{PBE}}
\end{displaymath}
这里$E_{x}^{\mathrn{HF}}$是\textrm{Fock}交换能(源自\textrm{Fock}交换势函数),而$E_x^{\mathrm{PBE}}$和$E_c^{\mathrm{PBE}}$是\textrm{PBE}的交换和相关能贡献。相比于其余杂化泛函，\textrm{PBE0}形式简洁，不含有任何拟合参数。

计算结构\textrm{POSCAR}文件

计算控制参数文件\textrm{INCAR}

\textrm{PBE0}计算的控制参数文件\textrm{PBE0/INCAR}

\textrm{POSCAR}采用的是金刚石结构的\textrm{Si}。

\textrm{INCAR}文件中指定了计算所使用的交换-相关泛函。%默认的\textrm{GGA}泛函无须指定。
设置计算控制参数\texttt{LHFCALC}来使用杂化泛函\textrm{PBE0}。而相关的计算控制参数\texttt{AEXX}，\texttt{AGGAC}和\texttt{ALDAC}将使用默认值。\footnote{因为\textrm{POTCAR}文件中已经确认原子数据集是采用\textrm{PBE}泛函计算得到的，因此原则上交换-相关泛函不必显式指定，而参数\texttt{AEXX}，\texttt{AGGAC}和\texttt{ALDAC}就是杂化泛函的系数。}

在\textrm{PBE}计算中，设置计算控制参数\texttt{ISMEAR}=-5，表示对\textrm{Fermi}能级的积分采用\textrm{Bl\"ochl}校正的四面体方法。可以这样得到最平滑的\textrm{DOS}曲线。取\texttt{LORBIT}=11，\textrm{VASP}会把投影的\textrm{DOS}写道\textrm{DOSCAR}和\textrm{vaspout.hf}文件中。

在\textrm{PBE0}计算中，电子弛豫采用阻尼算法，因此计算的能量要求必须是可变分的。直接优化算法比四面体方法更稳健。对于很多体系，如果采用杂化泛函，迭代收敛时对方法的稳健要求性更高。鉴于四面体方法的\textrm{Bl\'ochl}校正得到的能量是不可变分的，因此\textrm{Fermi}能级的积分采用\textrm{Gaussian}展宽方法，即设置计算控制参数\texttt{ISMEAR}=0，\texttt{SIGMA}=0.01。对阻尼算法来说，\textrm{Gaussian}展开得到的能量是可变分的。当体系达到收敛，再通过四面体方法(\textrm{ISMEAR}&=-5)完成非迭代(\texttt{ALGO}=\textrm{None})计算得到\textrm{DOS}和轨道投影(\texttt{LORBIT}=11)。

$\vec k$空间网格剖分是$\Gamma$点为中心的等空间网格。

通过该参数设置，带隙的计算精度预计会达到\textrm{0.1eV}以内。注意计算时，不要忘记对计算截断能和$\vec k$空间网格密度必须进行收敛测试。

执行\textrm{Bash}命令完成\textrm{VASP}计算

采用\textrm{Py4vasp}检查\textrm{PBE}计算的\textrm{DOS}曲线。

在\textrm{DOS}图中，带隙对应价带顶和导带底的能量差。在\textrmP{DOS}图中，\textrm{Fermi}能对应能量零点，带隙处没有任何电子态。

将\textrm{PBE}计算得到的\textrm{WAVECAR}复制到\textrm{PBE0}目录中，再次执行\textrm{Bash}命令完成\textrm{VASP}计算

%计算过程中，标准输出会在每个电子步迭代过程中，输出阻尼算法的额外的一些信息。

原则上，直接对比\textrm{PBE}和\textrm{PBE0}计算的总能是没有意义的，因为两者采用的是不同的交换-相关泛函。只有类似带隙这样的物理两才能对比。

通过\textrm{Python}脚本，对比两者计算的带隙。

\textrm{PBE0}计算得到的带隙(\textrm{1.84eV})比\textrm{PBE}的带隙(\textrm{0.62eV})明显大。在很多材料的电子计算中，和\textrm{GGA}泛函相比，杂化泛函得到更大的带隙，因为杂化泛函部分地校正了\textrm{GGA}泛函对带隙的低估。

为了直观地显示\textrm{PBE0}的带隙，还需要计算\textrm{DOS}。在\textrm{PBE0}计算的\textrm{INCAR}文件中，设置计算控制参数\texttt{ISMEAR}=-5，\texttt{ALGO}=\textrm{None}，\texttt{LORBIT}=11。此外还可以移除\texttt{SIGMA}和\texttt{TIME}参数。然后根据当前的\textrm{WAVECAR}，执行\textrm{Bash}命令，重启计算，并用\textrm{Py4vasp}模块绘制\textrm{DOS}曲线。

\textrm{PBE0}的\textrm{DOS}与\textrm{PBE}的不同，不只是带隙被扩大了，而且能带的相对位置和带宽都有差别。电子结构的这些微小改变，会对材料的性质产生很大的影响。

\subsection{\rm{PBE}、\rm{B3LYP}泛函和\rm{Hartree-Fock}计算的\ch{Ar}的带隙}
分别采用\textrm{PBE}、\textrm{B3LYP}泛函和\textrm{Hartree-Fock}方法计算\ch{Ar}的带隙。

\textrm{B3LYP}泛函最初是针对分子性质计算提出的半经验杂化泛函\textrm{(Semiempirical hybrid functional)}。所谓半经验，意味着其中有些参数是依据实验结果拟合的，如作为测试数据集的原子化能，电子-质子的亲和力，原子和分子的离子化能等。这些经验参数构成\textrm{B3LYP}泛函各组成部分的权重：
\begin{displaymath}
	\begin{aligned}
		E_{xc}^{\mathrm{B3LYP}}=&E_{x}^{\mathrm{B3LYP}}+E_{c}^{\mathrm{B3LYP}}\\
		E_{x}^{\mathrm{B3LYP}}=&0.8E_{x}^{\mathrm{LDA}}+E_{x}^{\mathrm{HF}}+0.72\Delta E_x^{\mathrm{B88}}\\
		E_{c}^{\mathrm{B3LYP}}=&0.19E_{c}^{\mathrm{VWN3}}+0.81E_{c}^{\mathrm{LYP}}
	\end{aligned}
\end{displaymath}
在\textrm{Hartree-Fock}计算中，交换项完全采用\textrm{Fock}交换形式(\texttt{AEXX}=1)，但完全不考虑相关能的贡献($E_c=0$)。

\textrm{B3LYP}和其余未屏蔽的泛函遗留了\textrm{Hartree-Fock}方法的一个致命缺陷，它没法在极限情况下满足均匀电子气模型。这个缺陷的根源是位于$\mathbf{q}=0$处的\textrm{Fock}势函数的\textrm{Fouriler}变换存在奇点\textrm{(singularity)}，即
\begin{displaymath}
	V_x^{\mathrm{HF}}(\mathbf{q})=-\dfrac{4\pi e^2}2\sum_{n\vec k}f_{n\vec k}\dfrac{\psi_{n\vec k}^{\ast}(\vec r^{\prime})\psi_{n\vec k}(\vec r)}{\mathbf{q}}
\end{displaymath}
即奇点的存在使\textrm{Fock}势函数完全压制点则的巡游属性。换言之，当处理具有显著巡游特征的体系，如金属和窄带半导体材料，不推荐\textrm{Hartree-Fock}方法和杂化泛函。\upciteP{Paier_2007}

本算例中，我们将计算\ch{Ar}的带隙，其带隙的实验值高达\textrm{14.4eV}。

计算的结构文件\textrm{POSCAR}

计算控制参数文件\textrm{INCAR}

\textrm{Hartree-Fock}方法的计算控制参数文件\textrm{HF/INCAR}

\textrm{B3LYP}泛函的计算控制参数文件\textrm{B3LYP/INCAR}

$\vec k$空间剖分网格文件\textrm{KPOINTS}

\textrm{POSCAR}文件给出的是\textrm{FCC}\ch{Ar}晶体

\textrm{INCAR}文件用于\textrm{PBE}计算

\textrm{HF/INCAR}，设置计算控制参数\texttt{LHFCAL}，开启\textrm{Fock}交换势贡献的计算，贡献权重由计算控制参数\texttt{AEXX}确定。因此这里当仅有\textrm{Fock}交换势贡献，完全没有半经验交换-相关泛函的贡献。因为\texttt{AEXX}=1，即使默认\texttt{GGA}=\textrm{PE}，其权重\texttt{AGGAX}=1-\texttt{AEXX}=0；而且\texttt{AGGAC}和\texttt{ALDAC}也将自动置0。

\textrm{B3LYP/INCAR}中，通过\textrm{GGA}泛函和各类交换-相关组分的系数的合理调节，可以设置多种半经验的泛函。

执行\textrm{Bash}命令完成\textrm{VASP}计算，并通过\textrm{Py4vasp}绘制\textrm{DOS}曲线。

将\textrm{PBE}计算得到的\textrm{WAVECAR}复制到\textrm{Hartree-Fock}计算和\textrm{B3LYP}计算目录下，并执行\textrm{Bash}命令完成\textrm{VASP}计算，通过\textrm{Py4vasp}得到三种方法的不同带隙，

计算结果显示，虽然\textrm{PBE}泛函得到的带隙是\textrm{8.5eV}，远小于实验值\textrm{14.4eV}，但是，\textrm{Hartree-Fock}方法的\textrm{17.6eV}和\textrm{B3LYP}泛函的\textrm{10.3eV}，两者的偏差也分别高达\textrm{3eV}以上。

\subsection{屏蔽杂化泛函计算\ch{MgO}的带隙}
为了得到与实验结果一致的\textrm{MgO}带隙，使用屏蔽杂化泛函并优化\textrm{Fock}势函数的系数。

屏蔽杂化泛函通过引入屏蔽参数$\mu$，将\textrm{Fock}势函数的长程衰减形式$1/r$拆解为短程部分\textrm{(short-range, SR)}和长程部分\textrm{(lon-range, LR)}两部分的贡献。
\begin{displaymath}
	E_{xc}^{\mathrm{screened}}=\dfrac14E_{x}^{\mathrm{HF,SR}}(\mu)+\dfrac34E_{x}^{\mathrm{PBE,SR}}(\mu)+E_{x}^{\mathrm{PBE,LR}}(\mu)+E_c^{\mathrm{PBE}}
\end{displaymath}
短程部分的极限回归到\textrm{PBE0}泛函形式。杂化泛函通过屏蔽参数阻断了长程部分的贡献，避免了\textrm{Fock}交换作用的奇点。而奇点的存在，是对各类导体属性压制的根源。此外，屏蔽杂化泛函通过误差函数，减少了计算开销。

固体物理中，常使用的泛函是\textrm{HSE06}。在该泛函中，优化的屏蔽参数取$\mu=0.2\mathrm{\AA}^{-1}$，对应的粒子间相互作用范围是$2/\mu\approx10\mathrm{\AA}$，基本上就是最近邻原子间相互作用的数倍。在\textrm{VASP}中，计算控制参数\texttt{HFSCREEN}来设置屏蔽长度。

带隙与选择的\textrm{Fock}交换势系数非常有关系。%这里采取实际的算例展示优化杂化系数，得到\textrm{MgO}的能带和实验结果接近的\textrm{7.7eV}。注意，这里牺牲了第一原理方法的预测能力。相反，该方法却又能更定量地描述模拟的物理。引用该泛函时，不要写成\textrm{HSE}，虽然文献里偶尔还这么写。
\textrm{HSE06}泛函表示\texttt{ASE}=0.25和\texttt{HFSCREEN}=0.2。

计算结构文件\textrm{POSCAR}

计算控制参数文件\textrm{INCAR}

$\vec k$空间剖分网格文件\textrm{KPOINTS}

与屏蔽杂化泛函有关的计算控制参数是\texttt{HFRCUT}和\texttt{HFAPHA}参数，分别表示屏蔽中的截断半径和衰减常数。

执行\textrm{Bash}命令执行\textrm{VASP}计算

并用\textrm{Py4vasp}绘制\textrm{DOS}曲线。

在屏蔽杂化目录下，执行\textrm{Bash}脚本修改杂化泛函中\textrm{Hartree-Fock}交换势系数\texttt{AEXX}，并完成\textrm{VASP}计算，

根据不同的\texttt{AEXX}的值，以实验的带隙\textrm{7.7eV}为参照，执行\textrm{Python}脚本，得到优化的\texttt{AEXX}的值。

优化的\texttt{AEXX}的值为$0.45\sim0.50$之间。

选择接近优化的\texttt{AEXX}的值，再改变屏蔽系数\texttt{HFSCREEN}，再进行计算，观察带隙是否对屏蔽参数的改变敏感。

短的屏蔽长度会增大\textrm{Fock}交换势的比例，进而影响带隙。

\subsection{采用\rm{PBE}和\rm{HSE06}泛函计算\ch{CaS}的能带}
对\ch{CaS}采用\textrm{HSE06}泛函，分别计算能带和带隙。

绘制能带图时，$\vec k$空间布点需要沿着第一\textrm{Brillouin}区的高对称性点的连线。由于对称性能意味着简并，所以这些高对称性的点和线特别需要注意。本算例中，确定最高占据带和最低未占据带的$\vec k$点(分别称为\textrm{HOMO}也称\textrm{LUMO})。如果两者对应同一$\vec k$点，则称为直接带隙，否则称为间接带隙。

确定绘制能带所需的高对称性$\vec k$点路径有两个软件，一个是\textrm{SeeK-path}软件，以\textrm{POSCAR}为输入，在\textrm{web}界面上选择高对称性$\vec k$点路径。%可以下载本算例的\textrm{POSCAR}，尝试能带路径\textrm{KPOINTS}文件。

\ch{CaS}的结构文件\textrm{POSCAR}

计算设置控制文件\textrm{INCAR}

\textrm{HSE06}计算控制参数文件\textrm{HSE06/INCAR}

\textrm{HSE06}能带计算控制参数文件\textrm{HSE06/band/INCAR}

$\vec k$空间网格剖分文件\textrm{KPOINTS}

能带绘制的高对称$\vec k$点路径\textrm{KPOINTS}

类似地，\textrm{HSE06}计算$\vec k$空间网格剖分文件\textrm{KPOINTS}

类似地，\textrm{HSE06}能带绘制高对称性$\vec k$点路径\textrm{KPOINTS}:~16个\textrm{IBZKPT}$\vec k$点和37个$\vec k$点绘制能带

\textrm{FCC}结构的\ch{CaS}的结构文件\textrm{POSCAR}

在\textrm{PBE}和\textrm{HSE06}泛函的自洽计算中，选用均匀$\vec k$空间网格剖分。\textrm{PBE}能带计算是线性格式，形成$\vec k$点路径。

杂化泛函的能带计算，既需要沿高对称性方向的$\vec k$点路径，还需要均匀网格剖分$\vec k$点来计算\textrm{Fock}交换势，因此\textrm{HSE06/band/KPOINTS}由16个\textrm{IBZ}$\vec k$点和37个高对称性$\vec k$点能带路径。

根据\textrm{Bash}命令完成\textrm{VASP}计算

进而完成\textrm{PBE}的能带计算

\textcolor{red}{采用\textrm{Py4vasp}模块绘制能带图}

在\textrm{HSE06}目录，通过\textrm{Bash}命令完成\textrm{VASP}计算

当\textrm{HSE06}计算过程中，对照\textrm{HSE06/band/KPOINTS}文件和\textrm{PBE}计算的\textrm{band/OUTCAR}文件中倒空间坐标。两者的$\vec k$空间能带路径相同，对于杂化泛函，为了计算\textrm{Fock}势，在能带计算过程中再次引入$\vec k$空间均匀布点。所以\textrm{KPOINTS}文件中，既有带权重的均匀$\vec k$空间网格剖分点，也有全部的绘制能带用的完整线性$\vec k$点。

通过\textrm{Bash}命令完成\textrm{HSE06}的能带计算

该计算会耗费一点时间，因为计算是由现有的\textrm{WAVECAR}文件开始的，体系总能已经达到收敛，因此\textrm{VASP}计算在达到控制参数\textrm{NELMIN}指定的迭代次数后就停止。为了得到体系的本征值，\textrm{Kohn-Sham}方程是在当前基组下对角化的，在额外的这些$\vec k$点上，波函数未经充分收敛，所以当前基组对能带计算的本征态未必足够。如果绘制的能带结构存在跳跃，为避免出现这种情况，至少应该让电子步在当前$\vec k$空间布点方案下至少迭代四次。

通过杂化泛函计算得到的能带图，包含均匀网格点和线性网格点，通过\textrm{Py4vasp}模块，得到两部分曲线的能带图。%采用活动模式缩放$L$点标记的能带部分。

仔细辨析导带，仍然存在微小的扭结。这种情况可以通过增加更多的迭代次数\textrm{(NELMIN)}或更多的能带\textrm{(NBANDS)}。

最后计算带隙，并通过\textrm{Py4vasp}绘制能带图。对比能带结构并得到带隙值。该带隙为间接带隙，因为\textrm{LUMO}位于$X$点，而\textrm{HOMO}位于$\Gamma$。
